% ---- pg-328 （书页 316） ----

其中
\[
\boldsymbol{\mu}=\left[\begin{array}{c}
\underset{(m\times1)}{\boldsymbol{\mu}_Y}\\
\hdashline
\underset{(r\times1)}{\boldsymbol{\mu}_Z}
\end{array}\right]
\qquad\text{及}\qquad
\boldsymbol{\Sigma}=\left[\begin{array}{c:c}
\underset{(m\times m)}{\boldsymbol{\Sigma}_{YY}}&\underset{(m\times r)}{\boldsymbol{\Sigma}_{YZ}}\\
\hdashline
\underset{(r\times m)}{\boldsymbol{\Sigma}_{ZY}}&\underset{(r\times r)}{\boldsymbol{\Sigma}_{ZZ}}
\end{array}\right]
\]
由结论 4.6, 当预测变量的值 $z_1,z_2,\cdots,z_r$ 给定时, $[Y_1,Y_2,\cdots,Y_m]'$ 的条件期望为
\phantomsection\label{eq:7-53}
\begin{equation}
E[\mathbf{Y}\mid z_1,z_2,\cdots,z_r]
=\boldsymbol{\mu}_Y+\boldsymbol{\Sigma}_{YZ}\boldsymbol{\Sigma}_{ZZ}^{-1}(\mathbf{z}-\boldsymbol{\mu}_Z)
\tag{7-53}
\end{equation}
该条件期望称为向量 $\mathbf{Y}$ 对 $\mathbf{Z}$ 的多元回归. 它由 $m$ 个一元回归构成. 例如, 条件均值向量的第一个分量为 $\mu_{Y_1}+\boldsymbol{\Sigma}_{Y_1Z}\boldsymbol{\Sigma}_{ZZ}^{-1}(\mathbf{z}-\boldsymbol{\mu}_Z)=E(Y_1\mid z_1,z_2,\cdots,z_r)$, 它使预测 $Y_1$ 的均方误差最小化. $m\times r$ 矩阵 $\boldsymbol{\beta}=\boldsymbol{\Sigma}_{YZ}\boldsymbol{\Sigma}_{ZZ}^{-1}$ 称为回归系数阵.

预测向量的误差
\[
\mathbf{Y}-\boldsymbol{\mu}_Y-\boldsymbol{\Sigma}_{YZ}\boldsymbol{\Sigma}_{ZZ}^{-1}(\mathbf{Z}-\boldsymbol{\mu}_Z)
\]
有平方与交叉乘积期望阵
\phantomsection\label{eq:7-54}
\begin{equation}
\begin{aligned}
\boldsymbol{\Sigma}_{YY\cdot Z}
&=E\left[\mathbf{Y}-\boldsymbol{\mu}_Y-\boldsymbol{\Sigma}_{YZ}\boldsymbol{\Sigma}_{ZZ}^{-1}(\mathbf{Z}-\boldsymbol{\mu}_Z)\right]
\left[\mathbf{Y}-\boldsymbol{\mu}_Y-\boldsymbol{\Sigma}_{YZ}\boldsymbol{\Sigma}_{ZZ}^{-1}(\mathbf{Z}-\boldsymbol{\mu}_Z)\right]'\\
&=\boldsymbol{\Sigma}_{YY}-\boldsymbol{\Sigma}_{YZ}\boldsymbol{\Sigma}_{ZZ}^{-1}(\boldsymbol{\Sigma}_{YZ})'
-\boldsymbol{\Sigma}_{YZ}\boldsymbol{\Sigma}_{ZZ}^{-1}\boldsymbol{\Sigma}_{ZY}
+\boldsymbol{\Sigma}_{YZ}\boldsymbol{\Sigma}_{ZZ}^{-1}\boldsymbol{\Sigma}_{ZZ}\boldsymbol{\Sigma}_{ZZ}^{-1}(\boldsymbol{\Sigma}_{YZ})'\\
&=\boldsymbol{\Sigma}_{YY}-\boldsymbol{\Sigma}_{YZ}\boldsymbol{\Sigma}_{ZZ}^{-1}\boldsymbol{\Sigma}_{ZY}
\end{aligned}
\tag{7-54}
\end{equation}
由于 $\boldsymbol{\mu}$ 和 $\boldsymbol{\Sigma}$ 一般是未知的, 它们必须用一个随机样本来估计, 以构造多元线性预测量并确定期望预测误差.

\mvthm{结论 7.14}{设 $\mathbf{Y}$ 和 $\mathbf{Z}$ 的分布为 $N_{m+r}(\boldsymbol{\mu},\boldsymbol{\Sigma})$, 此时 $\mathbf{Y}$ 对 $\mathbf{Z}$ 的回归为}
\[
\beta_0+\boldsymbol{\beta}\mathbf{z}
=\boldsymbol{\mu}_Y-\boldsymbol{\Sigma}_{YZ}\boldsymbol{\Sigma}_{ZZ}^{-1}\boldsymbol{\mu}_Z
+\boldsymbol{\Sigma}_{YZ}\boldsymbol{\Sigma}_{ZZ}^{-1}\mathbf{z}
=\boldsymbol{\mu}_Y+\boldsymbol{\Sigma}_{YZ}\boldsymbol{\Sigma}_{ZZ}^{-1}(\mathbf{z}-\boldsymbol{\mu}_Z)
\]
误差的平方与交叉乘积期望矩阵为
\[
E(\mathbf{Y}-\beta_0-\boldsymbol{\beta}\mathbf{Z})(\mathbf{Y}-\beta_0-\boldsymbol{\beta}\mathbf{Z})'
=\boldsymbol{\Sigma}_{YY\cdot Z}
=\boldsymbol{\Sigma}_{YY}-\boldsymbol{\Sigma}_{YZ}\boldsymbol{\Sigma}_{ZZ}^{-1}\boldsymbol{\Sigma}_{ZY}
\]
建立在一个容量为 $n$ 的随机样本基础上的回归函数的极大似然估计量为
\[
\hat{\beta}_0+\hat{\boldsymbol{\beta}}\mathbf{z}=\bar{\mathbf{Y}}+\mathbf{S}_{YZ}\mathbf{S}_{ZZ}^{-1}(\mathbf{z}-\bar{\mathbf{Z}})
\]
而 $\boldsymbol{\Sigma}_{YY\cdot Z}$ 的极大似然估计量则为
\[
\hat{\boldsymbol{\Sigma}}_{YY\cdot Z}
=\left(\frac{n-1}{n}\right)\left(\mathbf{S}_{YY}-\mathbf{S}_{YZ}\mathbf{S}_{ZZ}^{-1}\mathbf{S}_{ZY}\right)
\]

\medskip
\noindent\textbf{证明}\quad 回归函数和预测误差的协方差矩阵由结论 4.6 得到. 利用关系式
\[
\beta_0=\boldsymbol{\mu}_Y-\boldsymbol{\Sigma}_{YZ}\boldsymbol{\Sigma}_{ZZ}^{-1}\boldsymbol{\mu}_Z,\qquad
\boldsymbol{\beta}=\boldsymbol{\Sigma}_{YZ}\boldsymbol{\Sigma}_{ZZ}^{-1}
\]
\[
\beta_0+\boldsymbol{\beta}\mathbf{z}=\boldsymbol{\mu}_Y+\boldsymbol{\Sigma}_{YZ}\boldsymbol{\Sigma}_{ZZ}^{-1}(\mathbf{z}-\boldsymbol{\mu}_Z)
\]
\[
\boldsymbol{\Sigma}_{YY\cdot Z}=\boldsymbol{\Sigma}_{YY}-\boldsymbol{\Sigma}_{YZ}\boldsymbol{\Sigma}_{ZZ}^{-1}\boldsymbol{\Sigma}_{ZY}
=\boldsymbol{\Sigma}_{YY}-\boldsymbol{\beta}\boldsymbol{\Sigma}_{ZZ}\boldsymbol{\beta}'
\]
可由不变性[见\hyperref[eq:4-20]{式(4-20)}]得到上述量的极大似然估计量, 只需进行下列替代即可:
\[
\hat{\boldsymbol{\mu}}=\left[\begin{array}{c}\bar{\mathbf{Y}}\\ \hdashline \bar{\mathbf{Z}}\end{array}\right];
\qquad
\hat{\boldsymbol{\Sigma}}=\left[\begin{array}{c:c}\hat{\boldsymbol{\Sigma}}_{YY}&\hat{\boldsymbol{\Sigma}}_{YZ}\\ \hdashline \hat{\boldsymbol{\Sigma}}_{ZY}&\hat{\boldsymbol{\Sigma}}_{ZZ}\end{array}\right]
=\left(\frac{n-1}{n}\right)\mathbf{S}
=\left(\frac{n-1}{n}\right)\left[\begin{array}{c:c}\mathbf{S}_{YY}&\mathbf{S}_{YZ}\\ \hdashline \mathbf{S}_{ZY}&\mathbf{S}_{ZZ}\end{array}\right]
\]
\hfill$\blacksquare$

可以证明, $\boldsymbol{\Sigma}_{YY\cdot Z}$ 的无偏估计量为
\phantomsection\label{eq:7-55}
\begin{equation}
\left(\frac{n-1}{n-r-1}\right)\left(\mathbf{S}_{YY}-\mathbf{S}_{YZ}\mathbf{S}_{ZZ}^{-1}\mathbf{S}_{ZY}\right)
=\frac{1}{n-r-1}\sum_{j=1}^{n}(\mathbf{Y}_j-\hat{\beta}_0-\hat{\boldsymbol{\beta}}\mathbf{Z}_j)(\mathbf{Y}_j-\hat{\beta}_0-\hat{\boldsymbol{\beta}}\mathbf{Z}_j)'
\tag{7-55}
\end{equation}
